\documentclass[a4paper,12pt]{article}
\usepackage{fontspec}
\usepackage{geometry}
\geometry{margin=2cm}

\setmainfont{NotoSansLatIndic}[
  Path = ../../fonts/ttf-misc/NotoSansLatIndic/,
  Extension = .ttf,
  UprightFont = *-Regular,
  BoldFont = *-Bold,
]

\title{NotoSansLatIndic --- Merged Font Test}
\author{}
\date{}

\begin{document}
\maketitle

\section{Latin}
The quick brown fox jumps over the lazy dog. \\
\textbf{Bold: The quick brown fox jumps over the lazy dog.}

\section{Devanagari}
धर्मक्षेत्रे कुरुक्षेत्रे समवेता युयुत्सवः। \\
\textbf{Bold: धर्मक्षेत्रे कुरुक्षेत्रे समवेता युयुत्सवः।}

Conjunct test: क्ष त्र ज्ञ श्र द्ध क्र

\section{Tamil}
தர்மம் என்றால் கடமை. கர்மா என்பது செயல். \\
\textbf{Bold: தர்மம் என்றால் கடமை.}

\section{Telugu}
కర్మ అంటే చర్య. ధర్మం అంటే కర్తవ్యం. \\
\textbf{Bold: కర్మ అంటే చర్య.}

\section{Kannada}
ಕರ್ಮ ಎಂದರೆ ಕ್ರಿಯೆ. ಧರ್ಮ ಎಂದರೆ ಕರ್ತವ್ಯ. \\
\textbf{Bold: ಕರ್ಮ ಎಂದರೆ ಕ್ರಿಯೆ.}

\section{Bengali}
কর্ম মানে কাজ। ধর্ম মানে কর্তব্য। \\
\textbf{Bold: কর্ম মানে কাজ।}

\section{Malayalam}
കർമ്മം എന്നാൽ പ്രവൃത്തി. ധർമ്മം എന്നാൽ കർത്തവ്യം. \\
\textbf{Bold: കർമ്മം എന്നാൽ പ്രവൃത്തി.}

\section{Gujarati}
કર્મ એટલે ક્રિયા. ધર્મ એટલે ફરજ. \\
\textbf{Bold: કર્મ એટલે ક્રિયા.}

\section{Gurmukhi}
ਕਰਮ ਦਾ ਮਤਲਬ ਕੰਮ। ਧਰਮ ਦਾ ਮਤਲਬ ਫ਼ਰਜ਼। \\
\textbf{Bold: ਕਰਮ ਦਾ ਮਤਲਬ ਕੰਮ।}

\section{Oriya}
କର୍ମ ମାନେ କ୍ରିୟା। ଧର୍ମ ମାନେ କର୍ତ୍ତବ୍ୟ। \\
\textbf{Bold: କର୍ମ ମାନେ କ୍ରିୟା।}

\section{Mixed Script Paragraph}
The word \textbf{dharma} (धर्म / தர்மம் / ధర్మం / ಧರ್ಮ / ধর্ম / ധർമ്മം / ધર્મ / ਧਰਮ / ଧର୍ମ)
appears across all major Indian languages. The concept of कर्म (karma) is central to
Indian philosophy. In Tamil, it is written as கர்மா and in Telugu as కర్మ.

\section{Numbers}
Latin: 0 1 2 3 4 5 6 7 8 9 \\
Devanagari: ० १ २ ३ ४ ५ ६ ७ ८ ९ \\
Bengali: ০ ১ ২ ৩ ৪ ৫ ৬ ৭ ৮ ৯ \\
Tamil: ௦ ௧ ௨ ௩ ௪ ௫ ௬ ௭ ௮ ௯

\end{document}
